\chapter{Testing}
\label{ch:testing}

Kiss integrates JUnit 5\index{JUnit} for unit testing.  This chapter covers running the test suite, writing new tests, and testing strategies for Kiss applications.

\section{Running Tests}
\label{sec:test-running}

Run the full test suite with:

\begin{lstlisting}[style=shell]
./bld -v test
\end{lstlisting}

\noindent
This command:

\begin{enumerate}
  \item Compiles all source code (core, precompiled, and test classes).
  \item Packages the test classes with JUnit dependencies into \texttt{KissUnitTest.jar}.
  \item Executes the tests using JUnit Platform Console Launcher.
  \item Reports results to the terminal.
\end{enumerate}

\noindent
The \texttt{-v} flag enables verbose output, showing each test class and method as it runs.

\section{The Existing Test Suite}
\label{sec:test-existing}

Kiss ships with 19 test classes.  The core utility and database tests below live in \texttt{src/test/core/org/kissweb/} (with \texttt{QueryBuilderTest} under \texttt{database/}); an additional six OAuth test classes live under \texttt{oauth/} (\texttt{BearerTokenValidatorTest}, \texttt{JwksDiscoveryTest}, \texttt{oauth/as/AuthorizationServerTest}, and \texttt{oauth/client/\{OAuthClientConfigTest, OAuthClientStoreTest, PkceTest\}}):

\begin{longtable}{P{4.5cm}P{6cm}}
\toprule
\textbf{Test Class} & \textbf{What It Tests} \\
\midrule
\endhead
\texttt{ArrayUtilsTest} & Byte array primitive/boxed conversions. \\
\texttt{Base64Test} & Base64 encoding and decoding. \\
\texttt{DateTimeTest} & \texttt{DateTime} wrapper operations. \\
\texttt{DateUtilsTest} & YYYYMMDD date arithmetic and formatting. \\
\texttt{IniFileTest} & INI file parsing with sections and types. \\
\texttt{JsonPathTest} & Dot-path JSON navigation. \\
\texttt{NumberFormatTest} & Number formatting with masks. \\
\texttt{NumberUtilsTest} & Rounding and floating-point comparison. \\
\texttt{QueryBuilderTest} & SQL query building (the largest suite). \\
\texttt{StringUtilsTest} & String manipulation methods. \\
\texttt{TimeUtilsTest} & HHMM time operations. \\
\texttt{URLBuilderTest} & URL construction and parameter encoding. \\
\texttt{XMLTest} & XML parsing and transformation. \\
\bottomrule
\end{longtable}

\noindent
The \texttt{QueryBuilderTest} is the most comprehensive, containing over 100 tests that verify SELECT, WHERE, JOIN, subquery, CTE, and UNION generation against expected SQL output.

\section{Test Dependencies}
\label{sec:test-deps}

Kiss uses JUnit 5.11 (Jupiter).  The required JARs are included in the \texttt{libs/} directory:

\begin{itemize}
  \item \texttt{junit-jupiter-5.11.0.jar} --- aggregator module.
  \item \texttt{junit-jupiter-api-5.11.0.jar} --- annotations and assertions.
  \item \texttt{junit-jupiter-engine-5.11.0.jar} --- test engine.
  \item \texttt{junit-jupiter-params-5.11.0.jar} --- parameterized tests.
  \item \texttt{junit-platform-console-standalone-1.11.0.jar} --- command-line launcher.
\end{itemize}

\noindent
No external test framework installation is needed --- everything is bundled.

\section{Writing Tests}
\label{sec:test-writing}

Place test classes in \texttt{src/test/core/org/kissweb/} with the package declaration \texttt{org.kissweb}:

\begin{lstlisting}[style=java]
package org.kissweb;

import org.junit.jupiter.api.Test;
import static org.junit.jupiter.api.Assertions.*;

public class MyUtilTest {

    @Test
    public void testBasicOperation() {
        String result = MyUtil.process("input");
        assertEquals("expected", result);
    }

    @Test
    public void testNullHandling() {
        assertNull(MyUtil.process(null));
    }

    @Test
    public void testException() {
        assertThrows(IllegalArgumentException.class,
            () -> MyUtil.process(""));
    }
}
\end{lstlisting}

\subsection{Common Assertions}

\begin{longtable}{P{4.8cm}P{5.7cm}}
\toprule
\textbf{Assertion} & \textbf{Purpose} \\
\midrule
\endhead
\texttt{assertEquals(expected, actual)} & Verify equality. \\
\texttt{assertNotEquals(a, b)} & Verify inequality. \\
\texttt{assertTrue(condition)} & Verify condition is true. \\
\texttt{assertFalse(condition)} & Verify condition is false. \\
\texttt{assertNull(value)} & Verify value is null. \\
\texttt{assertNotNull(value)} & Verify value is not null. \\
\texttt{assertThrows(type, executable)} & Verify exception is thrown. \\
\texttt{assertDoesNotThrow(\allowbreak executable)} & Verify no exception. \\
\bottomrule
\end{longtable}

\subsection{Setup and Teardown}

Use \texttt{@BeforeEach} and \texttt{@AfterEach} for per-test initialization:

\begin{lstlisting}[style=java]
import org.junit.jupiter.api.BeforeEach;
import org.junit.jupiter.api.AfterEach;

public class ServiceTest {
    private Connection db;

    @BeforeEach
    void setUp() {
        // Initialize test state
    }

    @AfterEach
    void tearDown() {
        // Clean up resources
    }

    @Test
    void testService() {
        // test using initialized state
    }
}
\end{lstlisting}

\subsection{Parameterized Tests}

JUnit 5 supports parameterized tests for testing multiple inputs:

\begin{lstlisting}[style=java]
import org.junit.jupiter.params.ParameterizedTest;
import org.junit.jupiter.params.provider.CsvSource;

public class DateUtilsTest {

    @ParameterizedTest
    @CsvSource({
        "20260101, 2026, 1, 1",
        "20261231, 2026, 12, 31",
        "20260404, 2026, 4, 4"
    })
    void testDateComponents(int date,
            int year, int month, int day) {
        assertEquals(year, DateUtils.year(date));
        assertEquals(month, DateUtils.month(date));
        assertEquals(day, DateUtils.day(date));
    }
}
\end{lstlisting}

\section{Testing Database Operations}
\label{sec:test-db}

Database tests can use an in-memory SQLite database to avoid external dependencies.  Create a \texttt{Connection} with SQLite and a temporary file:

\begin{lstlisting}[style=java]
import org.kissweb.database.Connection;
import org.junit.jupiter.api.io.TempDir;
import java.nio.file.Path;

public class DatabaseTest {
    @TempDir
    Path tempDir;

    @Test
    void testInsertAndFetch() throws Exception {
        String dbPath = tempDir.resolve("test.db")
            .toString();
        Connection db = new Connection(
            Connection.ConnectionType.SQLite,
            dbPath);

        // Create table
        db.execute(
            "CREATE TABLE test_data (" +
            "id INTEGER PRIMARY KEY, " +
            "name TEXT, value REAL)");

        // Insert
        Record rec = db.newRecord("test_data");
        rec.set("name", "Alice");
        rec.set("value", 42.5);
        rec.addRecord();
        db.commit();

        // Fetch and verify
        Record fetched = db.fetchOne(
            "SELECT * FROM test_data " +
            "WHERE name = ?", "Alice");
        assertNotNull(fetched);
        assertEquals("Alice",
            fetched.getString("name"));
        assertEquals(42.5,
            fetched.getDouble("value"), 0.001);

        db.close();
    }
}
\end{lstlisting}

\begin{tipbox}
Use JUnit 5's \texttt{@TempDir} annotation to create temporary directories for SQLite databases.  The directory and its contents are automatically cleaned up after the test.
\end{tipbox}

\section{Testing Services}
\label{sec:test-services}

Service methods can be tested by constructing \texttt{JSONObject} inputs and calling the method directly.  Since service methods take standard parameters (\texttt{JSONObject}, \texttt{Connection}, \texttt{ProcessServlet}), you can create test instances:

\begin{lstlisting}[style=java]
import org.kissweb.json.JSONObject;

@Test
void testServiceMethod() throws Exception {
    // Set up database (SQLite for isolation)
    Connection db = createTestDatabase();
    setupTestData(db);

    // Build input
    JSONObject injson = new JSONObject();
    injson.put("employeeId", "EMP001");

    // Build output container
    JSONObject outjson = new JSONObject();

    // Call the service method directly
    EmployeeService svc = new EmployeeService();
    svc.getEmployee(injson, outjson, db, null);

    // Verify output
    assertEquals("Alice",
        outjson.getString("firstName"));
    assertEquals("Engineering",
        outjson.getString("department"));

    db.close();
}
\end{lstlisting}

\begin{notebox}
When testing services directly, pass \texttt{null} for the \texttt{ProcessServlet} parameter if the service does not use it.  If the service accesses session data via \texttt{servlet.getUserData()}, you will need to mock or stub the servlet.
\end{notebox}

\section{Testing Strategies}
\label{sec:test-strategies}

\begin{description}[style=nextline]
  \item[Unit tests] Test individual utility methods and pure logic.  These are fast and have no external dependencies.  Most of Kiss's existing tests are in this category.

  \item[Integration tests] Test service methods with a real (SQLite) database.  These verify that SQL queries, Record operations, and business logic work together correctly.

  \item[Frontend testing] Kiss's frontend is plain JavaScript without a build step.  Use browser developer tools for interactive debugging, or integrate a tool like Playwright or Cypress for automated UI testing.
\end{description}

\begin{warnbox}
Always use parameterized SQL queries in tests, just as in production code.  Never concatenate user input into SQL strings, even in test code, to maintain consistent security practices.
\end{warnbox}
